% DOCUMENT CLASS %
\documentclass[a4paper, 12pt]{article}   % for A4 Paper, 12pt fontsize.
\setcounter{tocdepth}{4}
\setcounter{secnumdepth}{3}

% PACKAGES%
\usepackage{amsfonts}   % For a basic mathfont (many will be replaced by stix)
\usepackage{amsmath}    % For basic math symbols
%\usepackage{bbm}        % For \mathbbm (better version of \mathbb)
\usepackage{mathrsfs}   % For \mathscr
\usepackage{hyperref}   % For footnotes
\usepackage{csquotes}   % For \textquote{}
\usepackage{IEEEtrantools}  % For better alignments
\usepackage{tikz}   % For a macro
\usepackage{listings} % For code
\usepackage{xcolor}   % For ... code
\usepackage[stretch=40]{microtype}
\usepackage{enumitem}
\usepackage{geometry}
\usepackage[color=gray]{todonotes}

% Language
\usepackage[english]{babel}

% If needed
\usepackage[nohyperlinks,smaller]{acronym}

% For tests
\usepackage{lipsum}

%%% ENVIRONMENTS %%%
\usepackage{amsthm}
\usepackage[nameinlink]{cleveref}

% Definition environment
\theoremstyle{definition}
\newtheorem{definition}{Definition}[section]
% Configure reference name
\crefname{definition}{Definition}{Definitions}

\theoremstyle{definition}
\newtheorem{exercise}{Exercise}[section]
\crefname{exercise}{Exercise}{Exercises}

\theoremstyle{definition}
\newtheorem{algorithm}{Algorithm}%[section]
\crefname{algorithm}{Algorithm}{Algorithms}

\theoremstyle{theorem}
\newtheorem{theorem}{Theorem}%[section]
\crefname{theorem}{Theorem}{Theorems}
%%%

%%% FONT CONFIGURATION %%%
\usepackage{fontspec}

% Updated version of Crimson Text
\setmainfont{Cochineal}
\setmonofont{Cascadia Code}[Scale=0.8]
\usepackage[cochineal]{newtxmath} % To use Cochineal as a Math Font
%%%

% MAKE DESCRIPTION ENVIRONMENTS HAVE ITALICIZED ITEMS %
\renewcommand{\descriptionlabel}[1]{\hspace{\labelsep}\textit{#1}}

% SETS %
\newcommand*{\R}{\mathbb R}    % Set of real numbers
\newcommand*{\N}{\mathbb N}    % Set of natural numbers
\newcommand*{\Z}{\mathbb Z}    % Set of integers
\newcommand*{\Q}{\mathbb Q}    % Set of rational numbers

% BASIC CONFIGURATION %
\pagestyle{plain}
\allowdisplaybreaks

% CODE %
\definecolor{codepurple}{HTML}{7f0055}
\definecolor{comment}{RGB}{0,128,0}
\definecolor{number}{RGB}{9,136,90}
\definecolor{string}{RGB}{163,21,21}

% Things that are the same in all styles
\lstdefinestyle{basic}{
    basicstyle=\ttfamily\footnotesize,
    breakatwhitespace=false,
    breaklines=true,
    frame=leftline,
    keepspaces=true,
    numbers=left,
    firstnumber=1,
    numbersep=5pt,
    showspaces=false,
    showstringspaces=false,
    showtabs=false,
    tabsize=2,
    commentstyle=\itshape\color{comment},
    stringstyle=\color{string},
    keywordstyle=\bfseries\color{codepurple},
    texcl=true,
    mathescape=true
}

\lstdefinestyle{default}{
    style=basic,
    language=c,
}

\lstset{style=default}

\title{Internet of Things and Wireless Networks}
\date{Summer Term 2026}
\author{Henri Heyden \\%
    \small stu240825 \\%
    \and%
    Parham Razaghian \\%
    \small stu258323%
}


\begin{document}
\maketitle

\section{Project Idea}
The idea we had for our IoT project is a BLE-based user device proximity/position detection.
Devices in a smart home may want to know if a user Bluetooth device is present nearby or in the same room without using a motion sensor.
This could be useful for triggering smart home devices or waking/suspending them from/to an energy saving mode, e.g., automatic lights and fans.
Furthermore, when using three or more devices, we can pinpoint the position of the user and use that data to wake specific services.

The latter point highlights another motivation---a privacy-based one.
We want to find out how realistic it can be for IoT devices to pinpoint the position and movement of Bluetooth devices.
Nefarious entities may have multiple motivations to find out this information.
It could be used to track user behaviour as it finds out which user is/was where and when.
This could be used for example to further personalize public advertisements based on people nearby.

Our approach is similar to the one, where a user device wants to find out its position, by scanning for signals of nearby devices whose position is known.
Our approach shows the other way, where these devices find out the position of the user device, without telling it.

\subsection{Concept}
We can implement the project idea using the following method.
\begin{enumerate}
    \item We place at least three devices in an arbitrary formation (optimally in a circle) around a user device with Bluetooth on\footnote{This is not restricted to smartphones, but also smartwatches, which often permanently have Bluetooth on.}.

    \item These devices continuously scan for foreign advertisements.

    \item Upon receiving an advertisement, we use the relative signal strength to heuristically calculate the distance to a device which is inside the polygon formed by our nodes.
        Furthermore, we can calculate the observed position in this area of the device directly by triangulation.
        For this we can use the signal strength in a weighted sum of the positions or solve a triangulation equation.

    \item Either way, that calculation may be too affording for simple IoT devices, which is why we want to send the measured information to a more powerful device.
\end{enumerate}

\subsection{Tasks}
Based on the preceding concept, we design three main tasks for us to solve.

\paragraph{Task 1: 5 Points}
Setup one device which relays the RSSIs of foreign advertisements to a devices which converts this measure to a distance.
Average values over time to reduce noise.

\paragraph{Task 2: 10 Points}
Setup three devices which relay the RSSIs to a smarter device which uses these signal strengths to triangulate the position of the device.
Again, reduce noise by averaging the samples and then output the position estimate.

\paragraph{Task 3: 15 Points}
Calculate a dynamic visualization on the smarter device to track the movement of the user device.
Also support multiple devices at the same time.
\end{document}
